% Shared preamble of the proof overview documents  <Folder>/<Name>_overview.tex.
% Each overview starts with
%     % Shared preamble of the proof overview documents  <Folder>/<Name>_overview.tex.
% Each overview starts with
%     % Shared preamble of the proof overview documents  <Folder>/<Name>_overview.tex.
% Each overview starts with
%     \input{../../scripts/overview_preamble.tex}
% (for ThinMatching/: \input{../scripts/overview_preamble.tex}) and is compiled from its own folder by
% scripts/build_overviews.py.
\documentclass[10pt,a4paper]{article}
\usepackage[T1]{fontenc}
\usepackage[utf8]{inputenc}
\usepackage[margin=2.2cm]{geometry}
\usepackage{amsmath,amssymb,amsthm}
\usepackage{array,booktabs,longtable}
\usepackage[dvipsnames]{xcolor}
\usepackage{enumitem}
\usepackage{tikz}
\usetikzlibrary{arrows.meta,positioning,shapes.multipart,fit,calc}
\usepackage[hidelinks]{hyperref}
\setlist{itemsep=2pt,topsep=4pt}
\renewcommand{\arraystretch}{1.2}

% Lean names, files and paper references
\newcommand{\lean}[1]{\textcolor{blue!50!black}{\texttt{\detokenize{#1}}}}
\newcommand{\file}[1]{\textcolor{OliveGreen!70!black}{\texttt{\detokenize{#1}}}}
\newcommand{\paperref}[1]{\textsf{#1}}
\newcommand{\hyp}[1]{\textcolor{BrickRed}{\textbf{#1}}}

% Boxes for the dependency diagrams
\tikzset{
  res/.style={draw, rounded corners, fill=blue!6, align=center, inner sep=4pt, font=\small},
  main/.style={res, fill=blue!18, very thick},
  cited/.style={res, fill=red!10, dashed},
  shared/.style={res, fill=gray!12},
  dep/.style={-{Stealth[length=2.2mm]}, thick},
}

\newcommand{\overviewtitle}[2]{%
  \begin{center}{\Large\bfseries #1}\\[4pt]{\small Proof overview of the Lean~4/Mathlib formalisation}\\[2pt]
  {\small #2}\end{center}}

% Standard paragraph describing the paper
\newcommand{\paperversion}{\noindent\textbf{Paper.} J.~Breitung, T.~Gervens, M.~Grohe.
  \emph{Graph Similarity Meets Optimal Transport}.\\ File \file{Graph_Similarity_Meets_Optimal_Transport.pdf} in the
  repository; statement, equation, page and line numbers refer to this PDF. \file{docs/NUMBERING.md} lists the Lean
  files and declarations of all numbered statements.\par\smallskip}

% (for ThinMatching/: \input{../scripts/overview_preamble.tex}) and is compiled from its own folder by
% scripts/build_overviews.py.
\documentclass[10pt,a4paper]{article}
\usepackage[T1]{fontenc}
\usepackage[utf8]{inputenc}
\usepackage[margin=2.2cm]{geometry}
\usepackage{amsmath,amssymb,amsthm}
\usepackage{array,booktabs,longtable}
\usepackage[dvipsnames]{xcolor}
\usepackage{enumitem}
\usepackage{tikz}
\usetikzlibrary{arrows.meta,positioning,shapes.multipart,fit,calc}
\usepackage[hidelinks]{hyperref}
\setlist{itemsep=2pt,topsep=4pt}
\renewcommand{\arraystretch}{1.2}

% Lean names, files and paper references
\newcommand{\lean}[1]{\textcolor{blue!50!black}{\texttt{\detokenize{#1}}}}
\newcommand{\file}[1]{\textcolor{OliveGreen!70!black}{\texttt{\detokenize{#1}}}}
\newcommand{\paperref}[1]{\textsf{#1}}
\newcommand{\hyp}[1]{\textcolor{BrickRed}{\textbf{#1}}}

% Boxes for the dependency diagrams
\tikzset{
  res/.style={draw, rounded corners, fill=blue!6, align=center, inner sep=4pt, font=\small},
  main/.style={res, fill=blue!18, very thick},
  cited/.style={res, fill=red!10, dashed},
  shared/.style={res, fill=gray!12},
  dep/.style={-{Stealth[length=2.2mm]}, thick},
}

\newcommand{\overviewtitle}[2]{%
  \begin{center}{\Large\bfseries #1}\\[4pt]{\small Proof overview of the Lean~4/Mathlib formalisation}\\[2pt]
  {\small #2}\end{center}}

% Standard paragraph describing the paper
\newcommand{\paperversion}{\noindent\textbf{Paper.} J.~Breitung, T.~Gervens, M.~Grohe.
  \emph{Graph Similarity Meets Optimal Transport}.\\ File \file{Graph_Similarity_Meets_Optimal_Transport.pdf} in the
  repository; statement, equation, page and line numbers refer to this PDF. \file{docs/NUMBERING.md} lists the Lean
  files and declarations of all numbered statements.\par\smallskip}

% (for ThinMatching/: \input{../scripts/overview_preamble.tex}) and is compiled from its own folder by
% scripts/build_overviews.py.
\documentclass[10pt,a4paper]{article}
\usepackage[T1]{fontenc}
\usepackage[utf8]{inputenc}
\usepackage[margin=2.2cm]{geometry}
\usepackage{amsmath,amssymb,amsthm}
\usepackage{array,booktabs,longtable}
\usepackage[dvipsnames]{xcolor}
\usepackage{enumitem}
\usepackage{tikz}
\usetikzlibrary{arrows.meta,positioning,shapes.multipart,fit,calc}
\usepackage[hidelinks]{hyperref}
\setlist{itemsep=2pt,topsep=4pt}
\renewcommand{\arraystretch}{1.2}

% Lean names, files and paper references
\newcommand{\lean}[1]{\textcolor{blue!50!black}{\texttt{\detokenize{#1}}}}
\newcommand{\file}[1]{\textcolor{OliveGreen!70!black}{\texttt{\detokenize{#1}}}}
\newcommand{\paperref}[1]{\textsf{#1}}
\newcommand{\hyp}[1]{\textcolor{BrickRed}{\textbf{#1}}}

% Boxes for the dependency diagrams
\tikzset{
  res/.style={draw, rounded corners, fill=blue!6, align=center, inner sep=4pt, font=\small},
  main/.style={res, fill=blue!18, very thick},
  cited/.style={res, fill=red!10, dashed},
  shared/.style={res, fill=gray!12},
  dep/.style={-{Stealth[length=2.2mm]}, thick},
}

\newcommand{\overviewtitle}[2]{%
  \begin{center}{\Large\bfseries #1}\\[4pt]{\small Proof overview of the Lean~4/Mathlib formalisation}\\[2pt]
  {\small #2}\end{center}}

% Standard paragraph describing the paper
\newcommand{\paperversion}{\noindent\textbf{Paper.} J.~Breitung, T.~Gervens, M.~Grohe.
  \emph{Graph Similarity Meets Optimal Transport}.\\ File \file{Graph_Similarity_Meets_Optimal_Transport.pdf} in the
  repository; statement, equation, page and line numbers refer to this PDF. \file{docs/NUMBERING.md} lists the Lean
  files and declarations of all numbered statements.\par\smallskip}

\begin{document}
\overviewtitle{Proposition 8: no linear bound for the local edit distance}{Folder \file{GsimOT/Proposition8_LocalEditLowerBound}}
\paperversion

\section{Statement}

\paragraph{Paper.} Proposition~8 (Section~3.2, p.~7, l.~262--263; construction and proof in Appendix~D.1,
pp.~25--28, l.~791--863): \emph{For every even $r\ge2$, there are graphs $G_r,H_r$ on $\Theta(r^3)$ vertices
such that $d_1(G_r,H_r)=r$ and $d^{\mathrm{OT}}_1(G_r,H_r)=1$.} Normalised, this is Theorem~3~(2) (p.~4,
l.~125): $\delta_1(G_r,H_r)=r\cdot\delta_1^{\mathrm{OT}}(G_r,H_r)$. Two remarks in the text concern the same
graphs: after Proposition~8 (p.~9, l.~295) ``Proposition~8 rules out any bound of the form
$\delta_1\le c\cdot\delta_1^{\mathrm{OT}}$'', and after Theorem~9 (p.~9, l.~318--320) ``there is no uniform
bound of the form $\delta_1\le C\cdot(\delta_1^{\mathrm{OT}})^\alpha$ with $\alpha>2/3$''.

\paragraph{Lean.} With $G_r$, $H_r$ the concrete graphs \lean{GsimOT.Prop8.Gr}, \lean{GsimOT.Prop8.Hr} on the
vertex type \lean{Vr} $r$ (core $\sqcup$ markers $\sqcup$ clique), and for every $r\in\mathbb N$ with
\hyp{$2\le r$ and $r$ even}:
\begin{itemize}
\item \lean{GsimOT.Prop8.proposition8}: $d_1(G_r,H_r)=r$, $d_1^{\mathrm{OT}}(G_r,H_r)=1$ and
  $r^3\le|V(G_r)|\le19r^3$; moreover \lean{card_Vr}: $|V(G_r)|=5r^3+10r^2-r+4$.
\item \lean{GsimOT.Prop8.theorem3_item2} (Theorem~3~(2)): $\delta_1(G_r,H_r)=r\cdot\delta_1^{\mathrm{OT}}(G_r,H_r)$.
\item \lean{GsimOT.Prop8.no_linear_bound}: for every $c\in\mathbb R$ some even $r\ge2$ has
  $\delta_1(G_r,H_r)>c\,\delta_1^{\mathrm{OT}}(G_r,H_r)$.
\item \lean{GsimOT.Prop8.no_power_bound}: for all $C\in\mathbb R$ and $\alpha>2/3$ some even $r\ge2$ has
  $\delta_1(G_r,H_r)>C\,(\delta_1^{\mathrm{OT}}(G_r,H_r))^\alpha$.
\end{itemize}
The Lean statements are exactly the paper's (with explicit constants for $\Theta(r^3)$).

\section{Where the proof lives}
\begin{longtable}{>{\raggedright\arraybackslash}p{5.2cm} >{\raggedright\arraybackslash}p{10.6cm}}
\toprule \textbf{File} & \textbf{Content} \\ \midrule\endhead \bottomrule\endfoot
\file{Construction.lean} & $\ell$, $n_X$, $n_M$, core $X$ and markers $M=C\cup T\cup B$; the adjacency table as
  patterns \lean{pat}; $G^0_r$, $H^0_r$ (\lean{G0}, \lean{H0}); degrees (\lean{deg_bip_le},
  \lean{deg_H0_marker}). Appendix~D.1, l.~792--809. \\
\file{Patterns.lean} & Consecutive patterns differ in one coordinate (l.~813--818, Fig.~6); how often each marker
  changes (\lean{nchg_*}); clock levels (\lean{clockLevel}) and their properties used in Lemma~28. \\
\file{LockedGraphs.lean} & Lemma~22 applied with markers sorted by degree, $\Delta=n_X$, $\tau=r$, $\delta=1$:
  \lean{Gr}, \lean{Hr}, the hypotheses \lean{lock_hyp_*}, $|V(G_r)|$ (\lean{card_Vr}); the shifts $\sigma_i$
  (\lean{sigma}). \\
\file{Lemma27_FractionalCost.lean} & Lemma~27, upper bound: $Q_r$ (\lean{Qr}, \lean{QrExt}) and
  $C_1(Q_r)\le1$ (\lean{costOne_Qr_le_one}); also $C_1(\sigma_i)\le r$ (\lean{costOneBij_sigmaExt_le}). \\
\file{Lemma28_BijectionCost.lean} & Lemma~28: the walk argument (\lean{walk_bits}, \lean{lemma28}). \\
\file{Proposition8.lean} & Lemma~27 ($=1$, \lean{lemma27}), Lemma~28 as stated (\lean{lemma28_dOne}),
  Proposition~8, Theorem~3~(2), the two remarks. \\
\file{Common/Observation21_DegreeBound.lean}, \file{Common/Lemma22_LockingMarkers.lean},
  \file{Common/Lemma29_Integrality.lean} & Observation~21, Lemma~22, Lemma~29 (shared). \\
\end{longtable}

\section{Proof layout}
\begin{center}
\begin{tikzpicture}
  \node[main] (t3) at (-4,0) {Theorem 3 (2)\\\lean{theorem3_item2}};
  \node[main] (rem) at (4,0) {remarks p.~9\\\lean{no_linear_bound}, \lean{no_power_bound}};
  \node[main] (p8) at (0,-1.8) {Proposition 8\\\lean{proposition8}};
  \node[res] (l27) at (-6.0,-3.9) {Lemma 27\\\lean{lemma27}};
  \node[res] (l28) at (2.0,-3.9) {Lemma 28\\\lean{lemma28}};
  \node[res] (card) at (6.4,-3.9) {$|V(G_r)|$\\\lean{card_Vr}};
  \node[res] (qr) at (-6.0,-6.3) {$C_1(Q_r)\le 1$\\\lean{costOne_Qr_le_one}};
  \node[res] (sig) at (-2.1,-6.3) {$C_1(\sigma_1)\le r$\\\lean{costOneBij_sigmaExt_le}};
  \node[res] (walk) at (1.9,-6.3) {walk on patterns\\\lean{walk_bits}};
  \node[shared] (o21) at (6.3,-6.3) {$L_\pi(v,\pi(v))\le C_1(\pi)$\\\lean{localCost_permMat_le_costOneBij}};
  \node[shared] (l29) at (-6.0,-8.8) {Lemma 29 (Common)\\\lean{lemma29_one_le_dOTOne}};
  \node[shared] (l22) at (-1.4,-8.8) {Lemma 22 (Common)\\\lean{lemma22_bij}, \lean{lemma22_cost}};
  \node[res] (pat) at (3.0,-8.8) {patterns, clock levels\\\file{Patterns.lean}};
  \draw[dep] (t3) -- (p8);
  \draw[dep] (rem) -- (p8);
  \draw[dep] (p8) -- (l27);
  \draw[dep] (p8) -- (l28);
  \draw[dep] (p8) -- (card);
  \draw[dep] (p8) -- (sig);
  \draw[dep] (l27) -- (qr);
  \draw[dep] (l27.west) to[out=200,in=160] (l29.west);
  \draw[dep, dashed] (l27) -- node[above, font=\scriptsize] {$G_r\not\cong H_r$} (l28);
  \draw[dep] (l28) -- (walk);
  \draw[dep] (l28) -- (o21);
  \draw[dep] (l28) -- (l22);
  \draw[dep] (qr) -- (l22);
  \draw[dep] (sig) -- (l22);
  \draw[dep] (walk) -- (pat);
\end{tikzpicture}
\end{center}
Arrows mean ``uses''. Blue: this folder; grey: \file{GsimOT/Common}. No cited result is used.
The dashed arrow is the derivation of $G_r\not\cong H_r$ from Lemma~28 (see Section~6). Lemma~27 and the
bound on $C_1(\sigma_1)$ also use the pattern facts of \file{Patterns.lean}. Observation~21 is used only inside
the proofs of Lemmas~22 and~29 in \file{GsimOT/Common} (see Section~6).

\section{The proof step by step}

\paragraph{Construction (Appendix~D.1, pp.~25--26, l.~792--809).} $\ell=r^2+2r$ (\lean{ell}),
$X=\{u\}\cup\{v_{i,k}\}$ (\lean{Core}), markers $c_{j,h}$, $t_{i,h}$, $b_i$ (\lean{Marker}; indices 0-based in
Lean). The table of l.~800 is the definition of \lean{pat} $r\,i\,t$, the pattern of branch $i$ at time $t$;
$G^0_r$, $H^0_r$ are \lean{G0}, \lean{H0}, and \lean{deg_H0_marker} is
$\deg_{H^0_r}(z)=\deg_{G^0_r}(z)+\mathbf 1[z\in C]$ (l.~802). Lemma~22 is applied with the markers sorted
non-decreasingly by degree (\lean{markerEnum}, \lean{markerEnum_mono}), $\Delta:=n_X$, $\tau:=r$, $\delta:=1$
(l.~803--809): \lean{Gr}, \lean{Hr}; the hypotheses of Lemma~22 are \lean{lock_hyp_deg_G},
\lean{lock_hyp_deg_H}, \lean{lock_hyp_delta}. In Lean, $G^0_r$ and $H^0_r$ are defined on the same type, so no
relabelling (l.~806--807) is needed.

\paragraph{Patterns (l.~810--818, Fig.~6).} $\chi_{G_r}(u)=\mathtt{pat}\,i\,0=0$ (\lean{pat_zero}),
$\chi_{H_r}(u)=\mathtt{pat}\,i\,\ell=\mathbf 1_C$ for even $r$ (\lean{pat_ell}); consecutive patterns differ in
at most one coordinate (\lean{ham_pat_succ_le_one}, via \lean{change_clock_iff}, \lean{change_tag_iff},
\lean{change_bit_iff}).

\paragraph{Lemma 27 (p.~27, l.~820--837).} The shifts $\sigma_i$ are \lean{sigma}/\lean{sigmaExt} and
$Q_r=\frac1r\sum_iP_{\sigma_i}$ is \lean{Qr}/\lean{QrExt}. Since $Q_r$ fixes markers and $K$, the second item of
Lemma~22 (\lean{lemma22_cost}) moves the computation to $(G^0_r,H^0_r)$. Supported core pairs cost their
Hamming distance $\le1$ (\lean{localCost_sigma_core}, \lean{ham_shift_le_one}); a marker costs the
$Q_r$-mass of the changes at it (\lean{localCost_sigma_marker}), bounded by $1$, $2/r$, $1$ for clocks, tags,
bits (\lean{nchg_clock_le}, \lean{nchg_tag_le}, \lean{nchg_bit_le}, \lean{sum_nchg_le}). Hence
$C_1(Q_r)\le1$ (\lean{costOne_Qr_base_le_one}, \lean{costOne_Qr_le_one}). The equality
$d_1^{\mathrm{OT}}=1$ (\lean{lemma27}) uses Lemma~29 (\lean{lemma29_one_le_dOTOne}) with $G_r\not\cong H_r$,
which is derived from Lemma~28 ($d_1\ge r>0$).

\paragraph{Lemma 28 (p.~27, l.~838--859).} In \lean{lemma28}: if $C_1(\pi)<r$, Lemma~22~(1)
(\lean{lemma22_bij} with $\tau=r$) shows that $\pi$ fixes all markers and permutes $X$; this gives the core
permutation $P$. The local cost at a core vertex is at least the Hamming distance of
$\chi_{G_r}(x)$ and $\chi_{H_r}(\pi(x))$, hence $<r$ (\lean{sum_Vr_ge_marker}). The walk
$u,P(u),\dots,P^s(u)=u$ ($s$ the least period) gives the pattern sequence $x_0,\dots,x_s$ (l.~842--844). The
rest is the purely combinatorial \lean{walk_bits}, which follows l.~844--857 line by line: clock levels
(\lean{clockLevel}) of consecutive patterns differ by $<r$ (\lean{clockLevel_le_add_ham}); $\beta$ is the first
index with full level, $\alpha<\beta$ the last with level $0$; all middle patterns lie on one branch $i$
because patterns of different branches with middle clock level differ in $\ge r$ tags
(\lean{pat_tag_of_middle}, \lean{ham_ge_of_middle_ne}); $\iota_j$ is the first index in $(\alpha,\beta]$ with
level $\ge jr$, $c(x_{\iota_j-1})\in((j-1)r,jr)$ forces block $j$ and bit value $j\bmod2$
(\lean{div_eq_of_clockLevel}, \lean{pat_bit_self}); counting changes (\lean{card_changes_step}) gives $\ge r$
alternations of $b_i$. Back in \lean{lemma28}, the $r$ alternations are $r$ distinct core vertices $x$ with
$A_{G_r}(b_i,x)\ne A_{H_r}(b_i,\pi(x))$, so $L_\pi(b_i,b_i)\ge r$, a contradiction.

\paragraph{Proof of Proposition 8 (pp.~27--28, l.~860--863).} $d_1^{\mathrm{OT}}=1$ is Lemma~27. For $d_1\le r$,
$\sigma_1$ fixes markers and $K$, so \lean{lemma22_cost} again reduces to $(G^0_r,H^0_r)$ and
\lean{costOne_sigma_base_le} shows that \emph{every} local cost of $\sigma_1$ is $\le r$
(\lean{costOneBij_sigmaExt_le}); with Lemma~28, \lean{dOne_Gr_Hr}: $d_1=r$. The order is \lean{card_Vr},
\lean{card_Vr_bounds}.

\paragraph{Theorem 3 (2) and the remarks (p.~4, l.~125; p.~9, l.~295 and l.~318--320).} \lean{theorem3_item2} divides by $n$.
\lean{no_linear_bound} takes $r>c$. \lean{no_power_bound}: $\delta_1=r/n$, $\delta_1^{\mathrm{OT}}=1/n$ and
$n\le19r^3$ give $n^{1-\alpha}\le19\,r^{3(1-\alpha)}$, so $\delta_1\le C(\delta_1^{\mathrm{OT}})^\alpha$ would imply
$r^{3\alpha-2}\le19C$, false for large $r$.

\section{Assumptions and trusted base}
No cited result and no hypothesis. All results depend only on \lean{propext}, \lean{Classical.choice},
\lean{Quot.sound} (checked by \file{scripts/CheckAxioms.lean}).

\section{Notes on the formalisation}
\begin{itemize}
\item \textbf{$G_r\not\cong H_r$.} The proof of Lemma~27 obtains $d_1^{\mathrm{OT}}(G_r,H_r)\ge1$ from Lemma~29
  since $G_r\not\cong H_r$ (p.~27, l.~837). The Lean derives $G_r\not\cong H_r$ from Lemma~28: every bijection
  has cost $\ge r>0$, while an isomorphism has cost $0$. It does not use the argument of the proof idea
  (p.~8, l.~282--283), which obtains the equality from Observation~21 since $u$ is isolated in $G_r$ but
  $H_r$ has no isolated vertex.
\item Lemma~27: the paper states the local costs at clock markers, tags and bits ($1$, $2/r$, $1$; l.~836--837);
  the Lean proves these values as upper bounds, which give $C_1(Q_r)\le1$.
\item Lemma~28: at a core vertex $x$ the Lean uses that $L_\pi(x,\pi(x))$ is at least the Hamming distance of
  the patterns (the paper: equal to it, l.~840--841). For patterns on different branches only Hamming
  distance $\ge r$ is used (the paper notes $\ge 2r$, l.~852).
\item The order $\Theta(r^3)$ is stated with explicit constants, $r^3\le|V(G_r)|\le19r^3$
  (\lean{card_Vr_bounds}), derived from the exact count \lean{card_Vr}.
\item The asymptotics $\delta_1^{\mathrm{OT}}=\Theta(r^{-3})$, $\delta_1=\Theta(r^{-2})$ (p.~9, l.~318--319) are
  not stated separately; their consequence (no exponent $\alpha>2/3$) is \lean{no_power_bound}.
\end{itemize}

\section{How to check this result}
\begin{itemize}
\item Build: \texttt{lake build GsimOT.Proposition8\_LocalEditLowerBound.Proposition8} (or \texttt{bash build.sh}
  for the whole library, which also runs \file{scripts/CheckAxioms.lean}).
\item Axioms: \texttt{\#print axioms GsimOT.Prop8.proposition8} (and \texttt{theorem3\_item2},
  \texttt{no\_linear\_bound}, \texttt{no\_power\_bound}); expected: \lean{propext}, \lean{Classical.choice},
  \lean{Quot.sound}.
\item \file{Proposition8_correspondence.pdf} in this folder lists every paper item with its Lean declaration.
\end{itemize}
\end{document}
